
\SGARefTarget{sga3:xxiv:1.5.2}
\subsubsection*{1.5.2. Existence.}
It suffices to prove existence locally for the \textup{(fpqc)}
topology. Let
\[
  (T,M,R,\Delta,(X_\alpha)_{\alpha\in\Delta})
  \quad\text{and}\quad
  (T',M',R',\Delta',(X'_{\alpha'})_{\alpha'\in\Delta'})
\]
be the two pinnings. By conjugacy of maximal tori, we may assume
\(T=T'\). After restricting \(S\), we may assume that the isomorphism
\[
  D_S(M)\simeq D_S(M')
\]
comes from an isomorphism \(M\simeq M'\) carrying \(R\) onto \(R'\);
we are thus reduced to the situation
\[
  T=T',\qquad M=M',\qquad R=R'.
\]
Since the systems of simple roots are conjugate under the Weyl group
\textup{(\SGARefLink{sga3:XXI:3.3.7}{Exp.~XXI, 3.3.7})}, we may also assume
\(\Delta=\Delta'\). For every \(\alpha\in\Delta\), there is then a
scalar \(z_\alpha\in\mathbb G_m(S)\) such that
\[
  X'_\alpha=z_\alpha X_\alpha,
\]
and it is enough to construct, locally for the \textup{(fpqc)}
topology, a section \(t\) of \(T\) such that
\(\alpha(t)=z_\alpha\) for every \(\alpha\in\Delta\). But the
morphism
\[
  T\longrightarrow(\mathbb G_{m,S})^\Delta
\]
whose components are the \(\alpha\), \(\alpha\in\Delta\), is dual to
an injection \(\mathbb Z^\Delta\to M\), and is therefore faithfully
flat. This completes the proof of \SGARefLink{sga3:xxiv:1.5.2}{1.5.2}, and hence of \SGARefLink{sga3:xxiv:1.3}{1.3}.

\SGARefTarget{sga3:xxiv:1.6}
\paragraph{\SGARefLink{sga3:xxiv:1.6}{Corollary 1.6}.}
\emph{Let \(S\) be a scheme and \(G\) a reductive \(S\)-group. The
following conditions are equivalent:}
\begin{enumerate}
\item[\SGAInlineRefTarget{sga3:xxiv:1.6:item:i}(i)] \(\underline{\Aut}_{S\text{-}\mathrm{gr.}}(G)\) is affine
  \emph{\textup{(respectively of finite type, respectively of finite
  presentation, respectively quasi-compact)}} over \(S\);
\item[\SGAInlineRefTarget{sga3:xxiv:1.6:item:ii}(ii)] \(\underline{\Autext}(G)\) is finite over \(S\);
\item[\SGAInlineRefTarget{sga3:xxiv:1.6:item:iii}(iii)] for every \(s\in S\),
  \[
    \operatorname{rg}_{\mathrm{red}}(G_s)
    -\operatorname{rg}_{\mathrm{ss}}(G_s)\leq 1.
  \]
\end{enumerate}

Indeed, since \(\operatorname{ad}(G)\) is affine, flat, and of finite
presentation over \(S\), the morphism
\[
  p:\underline{\Aut}_{S\text{-}\mathrm{gr.}}(G)
  \longrightarrow\underline{\Autext}(G)
\]
is affine, faithfully flat, and of finite presentation.

If \(\underline{\Autext}(G)\) is finite over \(S\), it is affine over
\(S\), and hence so is
\(\underline{\Aut}_{S\text{-}\mathrm{gr.}}(G)\); this proves
\textup{\SGARefLink{sga3:xxiv:1.6:item:ii}{(ii)}}\(\Rightarrow\)\textup{\SGARefLink{sga3:xxiv:1.6:item:i}{(i)}}. If
\(\underline{\Aut}_{S\text{-}\mathrm{gr.}}(G)\) is quasi-compact over
\(S\), it is of finite presentation over \(S\), since it is in any
case locally of finite presentation and separated over \(S\). By
\SGARefLink{sga3:v:proposition:9:1}{Exp.~V, 9.1}, \(\underline{\Autext}(G)\) is then of finite presentation
over \(S\), and hence finite; this proves
\textup{\SGARefLink{sga3:xxiv:1.6:item:i}{(i)}}\(\Rightarrow\)\textup{\SGARefLink{sga3:xxiv:1.6:item:ii}{(ii)}}. Finally, to prove the
equivalence of \textup{\SGARefLink{sga3:xxiv:1.6:item:ii}{(ii)}} and \textup{\SGARefLink{sga3:xxiv:1.6:item:iii}{(iii)}}, one may suppose that
\(G\) is split, and the assertion reduces to \SGARefLink{sga3:XXI:6.7.8}{Exp.~XXI, 6.7.8}.

\SGARefTarget{sga3:xxiv:1.7}
\paragraph{\SGARefLink{sga3:xxiv:1.7}{Corollary 1.7}.}
\emph{Let \(S\) be a scheme and \(G\) a reductive \(S\)-group. Then}
\[
  \underline{\Aut}_{S\text{-}\mathrm{gr.}}(G)^0
  \simeq\operatorname{ad}(G).
\]

\SGARefTarget{sga3:xxiv:1.8}
\paragraph{\SGARefLink{sga3:xxiv:1.8}{Corollary 1.8}.}
\emph{Let \(S\) be a scheme, \(G\) a reductive \(S\)-group, and
\(H\) an \(S\)-group scheme that is smooth, affine, and has connected
fibres over \(S\). Then the \(S\)-functor}
\[
  \underline{\operatorname{Isom}}_{S\text{-}\mathrm{gr.}}(G,H)
\]
\emph{is representable by a smooth separated \(S\)-scheme, which is
affine over \(S\) if \(G\) is semisimple.}

Indeed, let \(U\) be the set of points \(s\in S\) for which \(H_s\)
is reductive; this is an open subset \textup{(\SGARefLink{sga3:xix:corollary:2:6}{Exp.~XIX, 2.6})}. If
\(S'\) is an \(S\)-scheme, \(H_{S'}\) is reductive if and only if
\(S'\to S\) factors through \(U\). It follows that the canonical
morphism
\[
  \underline{\operatorname{Isom}}_{S\text{-}\mathrm{gr.}}(G,H)
  \longrightarrow S
\]
factors through \(U\). We may therefore assume \(S=U\), and we are
reduced to the following corollary.
